%macbook 
\loai{Lí thuyết}
\begin{vidu}%[3P5-6.1B][Loai1]
Trong thí nghiệm Y-âng về giao thoa ánh sáng, hai khe được chiếu sáng bởi nguồn phát đồng thời ba bức xạ đơn sắc: đỏ, lam, lục. Trong quang phổ bậc một, tính từ vân trung tâm ta sẽ quan sát thấy các vân sáng đơn sắc theo thứ tự
\choice
{đỏ, lam, lục}
{lục, lam, đỏ}
{lục, đỏ, lam}
{\True lam, lục, đỏ}
\loigiai{
Bức xạ đơn sắc có bước sóng càng nhỏ thì có khoảng vân càng nhỏ nên khoảng cách tới vân sáng trung tâm càng bé. Vậy tính từ vân trung tâm ra thì thấy: lam, lục, đỏ.
}
\end{vidu} \haidong{20}
\loai{Xác định bề rộng của quang phổ bậc 1}
\begin{vidu}%[3P5-6.1K][Loai2]
Trong thí nghiệm của Y-âng về giao thoa ánh sáng với nguồn ánh sáng trắng, hai khe hẹp cách nhau 0,5 mm, khoảng cách hai khe tới màn là 2 m. Khoảng cách giữa vân sáng bậc 1 của ánh sáng màu đỏ có bước sóng dài nhất ($\lambda_{\text{đ}}=0,75\ \mu m$) và vân sáng bậc 1 của ánh sáng màu tím có bước sóng ngắn nhất ($\lambda_t=0,4\ \mu m $) trên màn (gọi là bề rộng quang phổ bậc 1) là
\choice
{2,8 cm}
{2,8 mm}
{1,4 cm}
{\True 1,4 mm}
\loigiai{
\begin{itemize}
\item Ta có: $i_\text{đ}=\dfrac{\lambda_{\text{đ}}D}{a}=\dfrac{0,75. 2}{0,5}=3\ mm;\ i_t=\dfrac{\lambda_{t}D}{a}=\dfrac{0,4. 2}{0,5}=1,6\ mm $.
\item Suy ra bề rộng quang phổ bậc 1: $3-1,6 = 1,4\ mm $.
\end{itemize}
}
\end{vidu} \haidong{20}
\begin{vidu}%[3P5-6.1K][Loai2]
Trong thí nghiệm Y-âng về giao thoa ánh sáng, khoảng cách giữa hai khe 0,3 mm, khoảng cách từ mặt phẳng chứa hai khe đến màn quan sát 2 m. Hai khe được chiếu bằng ánh sáng trắng. Khoảng cách từ vân sáng bậc 1 màu đỏ (bước sóng $0,76\ \mu m$) đến vân sáng bậc 1 màu tím (bước sóng $0,4\ \mu m$) cùng phía so với vân trung tâm là
\choice
{1,8 mm}
{2,7 mm}
{1,5 mm}
{\True 2,4 mm}
\loigiai{ 
\begin{itemize}
\item Ta có: $\Delta_k=x_{d(k)}-x_{t(k)}=k\dfrac{\left({\lambda_\text{đ}-\lambda_t}\right)D}{a}=2,4.10^{-3}\ m$.
\item 
\item 
\end{itemize}
}
\end{vidu} \haidong{20}
\loai{Xác định bề rộng của quang phổ bậc k}
\begin{vidu}%[3P5-6.1G][Loai3]
Trong thí nghiệm Y-âng, hai khe được chiếu bằng ánh sáng trắng có có bước sóng $0,4\  \mu m\leq\lambda \leq 0,75\ \mu m$. Khoảng cách giữa hai khe là 0,3 mm, khoảng cách từ hai khe đến màn quan sát là 2 m. Khoảng cách giữa vân sáng bậc 3 màu đỏ và vân sáng bậc 3 màu tím ở cùng một bên so với vân trung tâm là
\choice
{$\Delta x=11\ mm$}
{\True $\Delta x=7\ mm$}
{$\Delta x=9\ mm$}
{$\Delta x=13\ mm$}
\loigiai{
\begin{itemize}
\item Khoảng cách giữa vân sáng bậc 3 màu đỏ và vân sáng bậc 3 màu tím ở cùng một bên so với vân trung tâm chính là bề rộng quang phổ bậc 3 của ánh sáng trắng.
\item Bề rộng quang phổ bậc 3 ứng với \bth{k=3}
\begin{align*}
\Delta x_3=3\dfrac{\pbrace{\lambda_{\ct{đỏ}}-\lambda_{\ct{tím}}}D}{a}=3.\dfrac{(0,75-0,40)2}{0,3}=7\ mm. 
\end{align*}
\end{itemize}
}
\end{vidu} \haidong{20}
\begin{vidu}%[3P5-6.1K][Loai3]
Trong thí nghiệm Y-âng về giao thoa ánh sáng, hai khe được chiếu bằng ánh sáng trắng có bước sóng \bth{0,4\ \mu m\leq\lambda \leq 0,75\ \mu m}. Khoảng cách giữa hai khe là 0,8\ mm, khoảng cách từ mặt phẳng chứa hai khe đến màn quan sát là 2\ m. Bề rộng quang phổ bậc 1 và bậc 2 quan sát được trên màn là
\choice
{\bth{\Delta x_1=0,32\ mm} và \bth{\Delta x_2=0,64\ mm}}
{\True \bth{\Delta x_1=0,875\ mm} và \bth{\Delta x_2=1,75\ mm}}
{\bth{\Delta x_1=2,32\ mm} và \bth{\Delta x_2=4,64\ mm}}
{\bth{\Delta x_1=0,32\ mm} và \bth{\Delta x_2=2,64\ mm}}
\loigiai{
Áp dụng công thức tính bề rộng quang phổ bậc k ta có
\begin{itemize}
\item Bề rộng quang phổ bậc 1 ứng với \bth{k=1}
\begin{align*}
\Delta x_1=\dfrac{\pbrace{\lambda_{\ct{đỏ}}-\lambda_{\ct{tím}}}D}{a}=\dfrac{(0,75-0,40)2}{0,8}=0,875\ mm.
\end{align*}
\item Bề rộng quang phổ bậc 2 ứng với \bth{k=2}
\begin{align*}
\Delta x_2=2\dfrac{\pbrace{\lambda_{\ct{đỏ}}-\lambda_{\ct{tím}}}D}{a}=2.\dfrac{(0,75-0,40)2}{0,8}=1,75\ mm.
\end{align*}
\end{itemize}
}
\end{vidu} \haidong{20}

\loai{Thay đổi cấu trúc hệ Y-âng}
\begin{vidu}%[3P5-6.1K][Loai4]
Trong thí nghiệm Y-âng, hai khe được chiếu bằng ánh sáng trắng có có bước sóng $0,4\ \mu m\leq\lambda \leq 0,75\ \mu m$. Bề rộng quang phổ bậc hai là 1,6\ mm. Khi dịch chuyển màn ra xa hai khe thêm 0,5\ m thì bề rộng quang phổ bậc hai bây giờ là 2\ mm. Khoảng cách giữa hai khe là
\choice
{1,125 mm}
{0,715 mm}
{\True 0,875 mm}
{1,300 mm}
\loigiai{
Áp dụng công thức tính bề rộng quang phổ bậc k ta có
\begin{itemize}
\item Bề rộng quang phổ bậc 2 trước khi dịch chuyển
\begin{equation}\rm
\Delta x_2=\dfrac{2D}{a}\pbrace{\lambda_{\ct{đỏ}}-\lambda_{\ct{tím}}}\tag{1}.
\end{equation}
\item Bề rộng quang phổ bậc 2 sau khi dịch chuyển
\begin{equation}\rm
\Delta x'_2=\dfrac{2(D+0,5)}{a}\pbrace{\lambda_{\ct{đỏ}}-\lambda_{\ct{tím}}}\tag{2}.
\end{equation}
\item Lấy \bth{\dfrac{(1)}{(2)}\Rightarrow \dfrac{\Delta x_2}{\Delta x'_2}=\dfrac{D}{D+0,5}=\dfrac{1,6}{2}=0,8\Rightarrow D=2\ m }.
\item Thay \bth{D=2\ m } vào (1) hoặc (2) ta có
\begin{align*}
a=\dfrac{2D}{\Delta x_2}\pbrace{\lambda_{\ct{đỏ}}-\lambda_{\ct{tím}}}=\dfrac{2.2}{1,6}(0,75-0,4)=0,875\ mm.
\end{align*}
\end{itemize}
}
\end{vidu} \haidong{20}

\begin{vidu}%[3P5-6.1K][Loai4]
Trong thí nghiệm của Y-âng về giao thoa ánh sáng với nguồn ánh sáng trắng, hai khe hẹp cách nhau 1mm. Xét về một phía vân trung tâm, khoảng cách giữa vân sáng bậc 1 của ánh sáng màu đỏ có bước sóng dài nhất ($\lambda_{\text{đ}}=0,76\ \mu m$) và vân sáng bậc 1 của ánh sáng màu tím có bước sóng ngắn nhất (${{\lambda}_{t}}=0,38\ \mu m $) trên màn (gọi là bề rộng quang phổ bậc 1) lúc đầu đo được là 0,38 mm. Khi dịch màn ra xa hai khe thêm một đoạn thì bề rộng quang phổ bậc 1 trên màn đo được là 0,57  mm. Màn đã dịch chuyển một đoạn bằng
\choice
{60 cm}
{\True 50 cm}
{55 cm}
{45 cm}
\loigiai{
\begin{itemize}
\item Bề rộng quang phổ bậc 1 được tính bằng: $L={i_d}-{i_t}=\dfrac{D}{a}({{\lambda}_{d}}-{{\lambda}_{t}})=0,38D=0,38\Rightarrow D=1\ m $.
\item Khi dịch màn $L'=0,57=0,38D'\Rightarrow D'=1,5\ m $.
\item Vậy màn đã dịch 50 cm.
\item 
\item 
\end{itemize}
}
\end{vidu} \haidong{20}
\begin{vidu}%[3P5-6.1K][Loai4]
Trong thí nghiệm của Y-âng về giao thoa ánh sáng với nguồn ánh sáng trắng. Xét về một phía vân trung tâm, khoảng cách giữa vân sáng bậc 1 của ánh sáng màu đỏ có bước sóng dài nhất ( $\lambda_{\text{đ}}=0,76\ \mu m $) và vân sáng bậc 1 của ánh sáng màu tím có bước sóng ngắn nhất ( ${{\lambda}_{t}}=0,40\ \mu m $) trên màn (gọi là bề rộng quang phổ bậc 1) lúc đầu đo được là 0,72 mm. Khi dịch màn ra xa hai khe thêm một đoạn 60 cm thì bề rộng quang phổ bậc 1 trên màn đo được là 0,90 mm. Khoảng cách giữa hai khe hẹp trong thí nghiệm là
\choice
{2 mm}
{1,5 mm}
{\True 1,2 mm}
{1 mm}
\loigiai{
\begin{itemize}
\item Bề rộng quang phổ bậc 1 được tính bằng: $L={i_d}-{i_t}=\dfrac{D}{a}({{\lambda}_{d}}-{{\lambda}_{t}})=0,38\dfrac{D}{a}=0,72\Rightarrow \dfrac{D}{a}=\dfrac{36}{19} $.
\item Khi dịch màn: $0,38\dfrac{D+0,6}{a}=0,9\Rightarrow \dfrac{D}{a}+\dfrac{0,6}{a}=\dfrac{45}{19}\Rightarrow a\approx 1,267\ mm $.
\item 
\item 
\item 
\end{itemize}
}
\end{vidu} \haidong{20}
\begin{vidu}%[3P5-6.2K][Loai1]
Trong thí nghiệm Y-âng về giao thoa ánh sáng trắng, bước sóng từ 400 nm đến 750 nm. Hai khe cách nhau 1,5\ mm và cách màn  giao thoa 1,2 m. Trên màn giao thoa, phần giao nhau giữa quang phổ bậc 1 và bậc 2 có bề rộng là
\choice
{\True Không tồn tại}
{0,04 mm}
{0,96 mm}
{0,24 mm}
\loigiai{
Theo bài ra ta có
\begin{itemize}
\item Vị trí vân đỏ bậc 1: \bth{x_{\ct{đỏ}1}=\dfrac{\lambda_{\ct{đ}}.D}{a}=\dfrac{0,75.1,2}{1,5}=0,6\ mm}.
\item Vị trí vân tím bậc 2: \bth{x_{\ct{tím}2}=\dfrac{\lambda_{\ct{t}}.D}{a}=2.\dfrac{0,4.1,2}{1,5}=0,64\ mm}.
\item Vì $x_{\ct{tím}2}>x_{\ct{đỏ}1}$ nên không có vùng phủ nhau giữa quang phổ bậc 1 và bậc 2.
\end{itemize}
}
\end{vidu} \haidong{20}
\begin{vidu}%[3P5-6.2K][Loai1]
Trong thí nghiệm của Y-âng về giao thoa ánh sáng với nguồn ánh sáng trắng, khoảng cách hai khe là 2 mm và khoảng cách hai khe tới màn là 2 m. Xét về một phía vân trung tâm, khoảng giữa vân sáng bậc n của ánh sáng màu đỏ có bước sóng dài nhất ($\lambda_{\text{đ}}=0,76\ \mu m $) và vân sáng bậc n của ánh sáng màu tím có bước sóng ngắn nhất (${{\lambda}_{t}}=0,38\ \mu m $) trên màn gọi là quang phổ bậc n. Vùng phủ nhau giữa quang phổ bậc hai và quang phổ bậc ba có bề rộng là
\choice
{\True 0,38 mm}
{1,14 mm}
{0,76 mm}
{1,52 mm}
\loigiai{
\begin{itemize}
\item Vùng phủ nhau giữa quang phổ bậc 2 và quang phổ bậc 3 chính là khoảng cách giữa vạch đỏ bậc 2 và vach tím bậc 3: $\Delta i=2{i_d}-3{i_t}=\dfrac{D}{a}(2{{\lambda}_{d}}-3{{\lambda}_{t}})=0,38\ mm $.
\item 
\item 
\item 
\end{itemize}
}
\end{vidu} \haidong{20}
\begin{vidu}%[3P5-6.2K][Loai1]
Trong thí nghiệm của Y-âng về giao thoa ánh sáng với nguồn ánh sáng trắng, khoảng cách hai khe là 0,76 mm và khoảng cách hai khe tới màn là 1,25 m. Xét về một phía vân trung tâm, khoảng giữa vân sáng bậc n của ánh sáng màu đỏ có bước sóng dài nhất ($\lambda_{\text{đ}}=0,76\ \mu m $) và vân sáng bậc n của ánh sáng màu tím có bước sóng ngắn nhất (${{\lambda}_{t}}=0,38\ \mu m $) trên màn gọi là quang phổ bậc n. Vùng phủ nhau giữa quang phổ bậc 4 và quang phổ bậc 5 có bề rộng
\choice
{1,25 mm}
{\True 1,875 mm}
{2,5 mm}
{3,75 mm}
\loigiai{
\begin{itemize}
\item Vùng phủ nhau giữa quang phổ bậc 4 và quang phổ bậc 5 chính là khoảng cách giữa vạch đỏ bậc 4 và vach tím bậc 5: $\Delta i=4{i_d}-5{i_t}=\dfrac{D}{a}(4{{\lambda}_{d}}-5{{\lambda}_{t}})=1,875\ mm $.
\item 
\item 
\item 
\end{itemize}
}
\end{vidu} \haidong{20}
